\chapter{Apolipoprotein B (ApoB)}

\section*{What Is This Test?}
Apolipoprotein B is the main structural protein of atherogenic lipoprotein particles, including LDL, VLDL remnants, IDL, and lipoprotein(a). ApoB estimates the number of circulating atherogenic particles and can add information when LDL cholesterol does not reflect particle burden.

\section*{Alternative Names}
ApoB; Apolipoprotein B-100; Atherogenic Particle Number.

\section*{Why This Test Is Ordered}
ApoB may be useful for cardiovascular-risk assessment when triglycerides are elevated, metabolic syndrome or diabetes is present, LDL cholesterol and non-HDL cholesterol are discordant, or residual risk remains uncertain despite standard lipid testing. It is not a replacement for clinical risk assessment or a routine stand-alone screening test for everyone.

\section*{How This Test Is Performed}
A blood sample is collected and ApoB is measured by an immunochemical assay. Fasting is often not required, although fasting may be requested when the test is obtained with triglycerides or a broader lipid panel.

\section*{Preparation, Risks, and Limitations}
Risks are limited to venipuncture. Acute illness, major changes in diet or weight, and lipid-lowering therapy can change results. Use the same laboratory when following trends, and interpret the result with LDL-C, non-HDL-C, triglycerides, diabetes status, blood pressure, and overall cardiovascular risk.

\section*{Understanding Results}
Higher ApoB indicates a greater number of atherogenic particles and generally higher atherosclerotic cardiovascular risk. Treatment thresholds are risk-specific rather than universal; a clinician may use ApoB to refine treatment intensity when standard cholesterol measures are ambiguous.

\section*{References}
American College of Cardiology and American Heart Association guidance on cholesterol and cardiovascular-risk assessment.

\clearpage
\section*{Understanding Results and Follow-up}
The laboratory report should identify the specimen, method, units, reference interval or decision limit, and whether the result is positive, negative, abnormal, or inconclusive. Reference intervals vary with age, sex, pregnancy, specimen type, assay, and laboratory; a result outside a range is not automatically a diagnosis, and a result within a range does not always exclude disease. Discuss unexpected results with the ordering clinician, who can decide whether repeat or confirmatory testing, treatment, or referral is appropriate. Do not change medicines or delay urgent care based on an isolated result.


\section*{Regulatory Status and Authorization Date}
Regulatory status applies to the specific assay, instrument, manufacturer, and intended use rather than to the test name alone. No single approval date can be assigned to this test category. For a commercial assay, verify the current product in the relevant regulator's database: the U.S. Food and Drug Administration (FDA) clearance, approval, or authorization; India's Central Drugs Standard Control Organisation (CDSCO) licence or permission; the European Union In Vitro Diagnostic Regulation (EU IVDR) CE certificate issued through the applicable conformity-assessment route; or the United Kingdom Medicines and Healthcare products Regulatory Agency (MHRA) registration and UKCA/CE status. Laboratory-developed tests may not have an individual FDA, CDSCO, EU IVDR, or MHRA approval date.
\clearpage
